\documentclass[10pt]{article}
\usepackage[margin=20mm]{geometry}
\usepackage{amsmath,amssymb,mathtools,fontspec,unicode-math,xurl,hyperref}
\IfFontExistsTF{DejaVu Serif}{\setmainfont{DejaVu Serif}}{\setmainfont{Latin Modern Roman}}
\setmathfont{latinmodern-math.otf}
\setlength{\parindent}{0pt}\setlength{\parskip}{6pt}\sloppy
\title{A technical review of
Central motives on parahoric flag varieties}
\author{AI-assisted technical review}
\date{27 October 2026}
\begin{document}
\maketitle
Abstract: The review covers the 58-page preprint arXiv:2403.11007v2, including its algebraic foundations, nearby-cycle constructions, geometric arguments, Wakimoto filtration and Hecke-algebra application. The specific augmentation bimodule used in Definition and Lemma 2.9 does not give a faithful derived tensor functor. An explicit nonzero map between finite free differential graded modules is sent to zero. This establishes a Major construction issue in the stated Koszul-duality step and the monodromy enhancement that uses it. It does not disprove the existence of a correctly implemented Koszul duality or the underlying geometric conclusions. A two-term complex identifies the missing action and explains why merely resolving the stated bimodule cannot repair the problem. Independent scientific adjudication remains pending.\par
\section*{Identity, version and scope}
The reviewed work is Central motives on parahoric flag varieties, by Robert Cass, Thibaud van den Hove and Jakob Scholbach [paper]. The exact source is arXiv:2403.11007v2, dated 6 December 2025, with 58 PDF pages. Page references below are one-based PDF pages and agree with the printed pagination. The associated journal identifier is DOI 10.4171/jems/1748; its public record dates online publication to 8 December 2025 [journal]. No equivalence between the journal typesetting and this preprint has been established. The frozen annual identity is M25-0364.\par
All 58 pages, including references, were read as extracted text. Fourteen selected source pages were also inspected as actual images: 1, 6–8, 12, 16–17, 22, 32, 38–40, 43 and 48. In particular, the source image on page 7 confirms the exact augmentation-bimodule construction and equations (2.7)–(2.9); the issue below is not an extraction artifact. The displayed diagrams on the inspected pages were checked against their surrounding arguments.\par
The assessment is Major, for the concrete foundational construction in Problem 1. This is an initial technical review, not independent acceptance or author adoption. External results on motivic six operations, geometric Satake, hyperbolic localization, local models and classical Hecke algebras were followed as cited dependencies; their complete proofs have not all been independently reconstructed.\par
\section*{Definitions, hypotheses and principal dependencies}
Equation reference: eq:E. Section 2.1.2 works in the category \(\mathrm{Ch}^{\mathbb Z}\) of complexes of internally \(\mathbb Z\)-graded abelian groups. Write \(\mathcal D^\mathbb Z\) for its derived category. The polynomial algebra \[\Sigma=\bigoplus_{n\geq0}\mathbb Z(n)=\mathbb Z[t]\] has \(t\) in cohomological degree zero and internal degree \(-1\), with primitive coproduct \(t\mapsto t\otimes1+1\otimes t\). The square-zero differential graded algebra on page 7 is \[
 I=\mathbb Z(-1)[-1],\qquad E=\mathbb Z\oplus I,\qquad I^2=0.\] In cohomological notation its generator \(\epsilon\) has degree \(1\) and internal degree \(1\); all differentials are zero. The source uses the right-induced projective model structure on \(E\)-modules, and a left-induced model structure on \(\Sigma\)-modules whose weak equivalences are the underlying quasi-isomorphisms. Both present the respective usual derived module categories.\par
The general motivic base \(S\) is connected and smooth of finite type over a Dedekind ring or a field (Notation 2.11). The geometric construction uses a split reductive group and facets in the closure of a fixed alcove. Assertions involving the mixed Tate heart retain the stated Beilinson–Soulé hypothesis, except where reduced motives remove that requirement. Integral motives and rationally defined monodromy are distinguished throughout the review.\par
The dependency relevant here is specific. Corollary 2.20, whose proof continues on page 12, invokes Lemma 2.9 to identify the reduced Tate category on \(\mathbb G_m\) with \(\operatorname{Nilp}\), and uses rational coefficients for the corresponding unreduced identification. Definition 2.24 then forms the functor with values in \(\operatorname{Nilp}_{\mathbb Q}\operatorname{DTM}(S)\). In the proof of Theorem 3.2, page 16 explicitly forms extension of scalars through this functor. The same foundational construction is repeated in several variables on page 22, and enters the monodromy-compatible convolution construction in Lemma 4.20, page 32. Thus this is an operative step in the paper’s enhanced nearby-cycle construction.\par
\section*{Substantive finding}
Problem 1 (The specified augmentation functor kills a nonzero morphism).  On page 7, the paragraph before (2.7) gives \(\mathbb Z\) its \(\Sigma\)–\(E\) bimodule structure by the two augmentations, in \(\mathrm{Ch}^{\mathbb Z}\). The derived functor of the displayed left Quillen functor \[\mathbb Z\otimes_E-:\operatorname{Mod}_E(\mathrm{Ch}^{\mathbb Z})^{\mathrm r}
       \longrightarrow\operatorname{Mod}_\Sigma(\mathrm{Ch}^{\mathbb Z})^{\mathrm l}\] is not faithful. Consequently the full-faithfulness assertion of Definition and Lemma 2.9, as constructed in (2.7)–(2.9), is false.\par
\subsection*{Complete admissible input}
Equation reference: eq:input. Use exactly the algebras in [eq:E] (\#eq:E), over \(\mathbb Z\), and the strict bimodule specified in the paper. Consider the left \(E\)-modules \[
 A=E\otimes_{\mathbb Z} I=E(-1)[-1],\qquad B=E.\] They are free \(E\)-modules on internally graded complexes consisting of one copy of \(\mathbb Z\) in a single degree. They are therefore cofibrant in the source model structure specified on page 7. They are also compact objects of the derived category. No boundedness, coefficient, finiteness or nilpotence restriction in Lemma 2.9 excludes these inputs.\par
Equation reference: eq:f. Let \(a=1\otimes\epsilon\) denote the free generator of \(A\), and let \(b=\epsilon a\). Define \[
 f:A\longrightarrow B,\qquad f(a)=\epsilon,\qquad f(b)=0.\] The map has cohomological degree zero and internal degree zero. Indeed, the bidegrees of \(a,b\) are \((1,1),(2,2)\), and those of \(1,\epsilon\) in \(B\) are \((0,0),(1,1)\). The only nonzero image has matching bidegree. It is \(E\)-linear because \(f(\epsilon a)=0=\epsilon f(a)\), and it is a chain map since both differentials vanish.\par
\subsection*{The morphism is nonzero in the derived category}
Equation reference: eq:nonzero. Since \(A\) is free on \(I\), the derived mapping complex is obtained from the ordinary graded mapping complex: \[\operatorname{RHom}_E(E\otimes I,E)\simeq \operatorname{RHom}_{\mathrm{Ch}^{\mathbb Z}}(I,E).\] The degree-zero, internal-degree-zero maps form \(\mathbb Z\), generated by the inclusion of \(I\) into \(E\). All differentials in this mapping complex are zero, so the indicated class is not a boundary. Equivalently, a chain homotopy would contribute \(d_Bh+hd_A=0\) and cannot kill \(f\). Thus \[
 0\neq[f]\in\operatorname{Hom}_{\operatorname{Mod}_E(\mathcal D^\mathbb Z)}(A,B).\] This argument uses cofibrant inputs; it does not confuse an arbitrary underived map with a nonzero derived morphism.\par
\subsection*{Its image under the stated derived functor is zero}
Equation reference: eq:zero. Let \(F\) be the derived functor in (2.8), obtained from (2.7) as in the source proof. On the two cofibrant objects [eq:input] (\#eq:input) and their map, its value is computed by the stated ordinary tensor construction: \[F(A)=\mathbb Z\otimes_E(E\otimes I)\cong I,\qquad
 F(B)=\mathbb Z\otimes_E E\cong\mathbb Z.\] The left \(\Sigma\)-action on both is the augmentation action. The right \(E\)-action on the bimodule \(\mathbb Z\) also is the augmentation action, so \[
 F(f)(1\otimes a)=1\otimes\epsilon
       =(1\cdot\epsilon)\otimes1=0.\] It follows that \(F(f)=0\) as an actual chain map, and hence as a derived \(\Sigma\)-module morphism. Equations [eq:nonzero] (\#eq:nonzero) and [eq:zero] (\#eq:zero) contradict faithfulness. The same calculation works after tensoring all data with \(\mathbb Q\); changing to rational coefficients does not remove this defect.\par
An independent exact calculation retained with this review verifies the square-zero algebra, all bidegrees, \(E\)-linearity of the matrix for \(f\), the rank-one degree-zero mapping group, and the zero map induced in the augmentation quotient. This is a finite corroboration of the explicit witness. The passage to the derived category is justified by the freeness argument above, not by numerical testing.\par
\subsection*{Why the proof’s endomorphism calculation does not settle the matter}
The resolution \[0\longrightarrow\Sigma(1)\xrightarrow{\,t\,}\Sigma\longrightarrow\mathbb Z\longrightarrow0\] does compute the graded groups of \(\operatorname{RHom}_\Sigma(\mathbb Z,\mathbb Z)\) and proves the needed compactness of \(\mathbb Z\). It does not show that the particular adjunction unit supplied by the augmentation bimodule is a quasi-isomorphism. An abstract identification of the derived endomorphism algebra with a square-zero algebra cannot replace verification of that unit map. The morphism [eq:f] (\#eq:f) tests its positive-degree information directly and is killed by the displayed functor.\par
One might instead read the notation \(\mathbb Z\) as a derived Koszul bimodule carrying a nontrivial higher action. Such a construction can have different behavior. However, the source explicitly chooses the two augmentation actions in the ordinary category of complexes, then derives the resulting Quillen functor. That explicit construction fixes the interpretation tested here.\par
\subsection*{Repair boundary}
Equation reference: eq:Kaction. A two-term complex makes the missing information concrete. Set \[K^{-1}=\Sigma u\cong\Sigma(1),\qquad K^0=\Sigma,\qquad d(u)=t,\] where the notation means that \(u\) has cohomological degree \(-1\) and internal degree \(-1\). Give \(K\) the right \(E\)-action \[
 u\epsilon=1,\qquad 1\epsilon=0,\] extended \(\Sigma\)-linearly. Its bidegrees match, its square is zero, and the differential is compatible with the right action. As a left \(\Sigma\)-complex, \(K\) resolves \(\mathbb Z\). Nevertheless, the usual augmentation \(K\to\mathbb Z\) is not a map to the strict augmentation \(\Sigma\)–\(E\) bimodule: it sends \(u\epsilon\) to \(1\), whereas the augmentation of \(u\) multiplied by \(\epsilon\) is zero.\par
Thus the missing action cannot be supplied by merely taking a resolution within the derived category of the already specified bimodule. A corrected construction must specify a different derived bimodule or an equivalent Koszul-duality implementation, and identify its endomorphism action. The calculation [eq:Kaction] (\#eq:Kaction) isolates that necessary change; this review does not assert a full repair of the symmetric monoidal structure, its coherences, or all of the monodromy and convolution compatibilities.\par
\subsection*{Impact and severity}
This is a Major issue because an asserted foundational fully faithful functor is not even faithful, and that specific functor supplies the paper’s description of \(\operatorname{Nilp}\) and the subsequent monodromy enhancement. The explicit uses in Corollary 2.20, Definition 2.24, the proof of Theorem 3.2, and Section 3.10 must be revisited with a correct construction. The tensor identifications in the monodromy triangle, Proposition 3.9, also require the correct module actions, rather than only an equality of underlying complexes.\par
The counterexample does not refute a properly implemented Koszul duality. Nor does it establish that the underlying central-motive, Whitney–Tate, Wakimoto or Hecke-algebra conclusions are false. In particular, the construction before retaining monodromy and the geometric arguments have substantial content independent of the specific failed action. The scope of the finding is the construction and its stated proof chain; a repaired proof could preserve the intended conclusions. No author correction or adoption of this repair has been verified.\par
\section*{Coverage of the remaining arguments}
The higher-algebra preliminaries were read with the compact-generation and rigidity assumptions retained. The distinctions between integral and rational formality of the motive of \(\mathbb G_m\), reduced motives, and the finite-characteristic obstruction discussed in Remark 2.19 were kept in scope. The problem established above concerns the later use of the specified bimodule, rather than the mere existence of the polynomial and square-zero algebras.\par
The nearby-cycle discussion was followed through its proper and smooth functoriality, projection-formula qualifications for trivial families, hyperbolic localization under the stated monodromicity and local linearizability hypotheses, realization comparisons, and the multivariable diagonal and iterated comparisons. Natural comparison maps were not treated as isomorphisms where the paper does not assert that.\par
In Section 4 the review followed the global and local convolution diagrams, the distinction between étale and Zariski local sections, finite-jet reductions, fixed-point and attractor descriptions, Demazure induction, and the noetherian induction detecting Tateness and anti-effectivity. The proof of Theorem 4.36 reduces the relevant assertions to the trivial family using hyperbolic localization and conservativity. The centrality and monoidality assertions in Theorems 4.38 and 4.41 were checked with their stated homotopy-category coherence scope. The mixed-Tate tensor-product assumption in Theorem 4.49 was retained; no stronger general centrality theorem is inferred.\par
Section 5 was read through standard and costandard convolution, the invertibility twist, the group-completion construction of Wakimoto functors, the bounded-support criterion and filtration, constant terms, and the shift and twist in Theorem 5.24. The passage to general parahorics in Theorem 5.30 uses a suitable Borel, smooth conservative pullback, and the two one-sided exactness arguments; it does not assume nearby cycles are generally t-exact at the outset.\par
Section 6 was checked for the sign in the standard basis, the quadratic relation, specialization to a finite-field Hecke algebra, the extra twist in the Bernstein elements, and the triangular basis argument in Theorem 6.9. The central map is asserted to be an isomorphism at Iwahori level; the review does not silently extend that conclusion to arbitrary parahorics. No additional substantive error is established by these checks. This statement records the actual review outcome and does not certify every imported theorem or every higher categorical coherence independently.\par
A bounded search of the exact arXiv record, journal record, public author pages, and title/code/correction queries did not locate accompanying author code, supplementary data or an applicable erratum. This does not establish their absence or the priority of the present observation. No author code was executed, and lack of code is not a mathematical defect.\par
\section*{Conclusion}
Item | Recorded result\par
Problem 1 | Major: explicit nonzero derived morphism killed by the specified functor.\par
Repair | A different bimodule action is necessary; full monoidal and geometric repair remains unverified.\par
Other conclusions | No independent counterexample to the geometric conclusions established.\par
Scientific decision | Pending; this is an initial review.\par
Human reviews | Empty.\par
Author adoption | Unverified.\par
\section*{Supporting checks}
The finite matrix witness uses the ordered bases \((a,\epsilon a)\) and \((1,\epsilon)\): \[R_\epsilon=\begin{pmatrix}0&0\\1&0\end{pmatrix},
 \qquad f=\begin{pmatrix}0&0\\1&0\end{pmatrix},
 \qquad (1\ \ 0)\,f\,\begin{pmatrix}1\\0\end{pmatrix}=0.\] For the illustrative complex \(K\), in the ordered \(\Sigma\)-basis \((1,u)\), the differential and right action are \[d_K=\begin{pmatrix}0&t\\0&0\end{pmatrix},
 \qquad r_\epsilon=\begin{pmatrix}0&1\\0&0\end{pmatrix}.\] Both squares and both mixed products vanish. The augmentation row times \(r_\epsilon\) is \((0\ \ 1)\), proving the failure of bimodule linearity stated above. The successful execution used the existing Python runtime and SymPy 1.13.3. An earlier execution reached the mathematical assertions but failed when serializing SymPy integers to JSON; its unchanged script, exit-one record and error output are preserved. That reporting failure is not attributed to the authors.\par
d7866edf012b3ab33150a549b81ca89dac43240dd1930b0564c7ac312a7f09b8.\par
The source manifest, metadata, complete extracted text, selected source images, actual visual observations, evidence and complete support inventory are retained in the allocated staging directories. Source image inspection covers the selected fourteen pages; full text reading covers all 58.\par
Compilation uses the existing approved Tectonic 0.17.0 runtime with shell escape disabled. Exact TeX, compiler and dependency payloads before and after compilation, commands, exit codes, PDF parsing, and actual final-page inspection are recorded separately. Rendering alone is not visual inspection, and delivery is marked complete only after the final PDF and those records are checked.\par
All substantive analysis and visual judgments in this report belong to the current scientific worker conversation. Native tools performed the explicitly logged arithmetic, byte operations, compilation and rendering. Main’s role is UI/logistics and byte transport. No human scientific review or independent acceptance is claimed.\par
9 R. Cass, T. van den Hove and J. Scholbach, Central motives on parahoric flag varieties, arXiv:2403.11007v2, 6 December 2025, 58 pages. https://arxiv.org/abs/2403.11007v2. EMS Press, Central motives on parahoric flag varieties, journal metadata, DOI 10.4171/jems/1748. https://ems.press/journals/jems/articles/14299339. Consulted for identity and publication metadata; the journal full text was not used as a substitute for the bound preprint.\par
\end{document}
